\section{Polynomial bases, shift equivariance, and Volterra structure}
\label{app:offsets}

Polynomial position weighting does not single out one coordinate basis on the
space of polynomials of degree below \(P\). Separating the subspace of position
summaries retained by the lift from the choice of basis clarifies three points
needed later: why \(P\) copies interpolate between unweighted signatures and full
sequence reconstruction, why positive-diagonal QR orthonormalization produces a
canonical discrete orthogonal basis for HAR and DSADS, and why the binomial
Pascal basis itself admits an exact triangular shift law and a position-only
Volterra representation.
We record standard polynomial basis and translation identities in the notation
needed for these coded signatures.

\subsection{From exact slot labels to polynomial moments}
An uncompressed representation of discrete positions would keep every time step
in its own dedicated block of \(C\) channels. Let \(e_0, \ldots, e_{R-1}\) denote
the standard basis of \(\R^R\). Assigning \(e_r \otimes z_r \in \R^R \otimes \R^C \cong \R^{RC}\)
to step \(r\) records both the increment \(z_r\) and its exact index \(r\), at
the cost of \(RC\) channels. Any linear code \(\mathcal A : \R^R \to \R^P\)
compresses these \(R\)-dimensional slot indicators into \(P\)-dimensional code
vectors \(\mathcal A e_r = b(r) \in \R^P\):
\begin{equation}\label{eq:offset-code}
 \mathcal A e_r = b(r), \qquad
 (\mathcal A \otimes I_C)(e_r \otimes z_r) = b(r) \otimes z_r = v_r,
\end{equation}
where we order coordinates copy-major so \(b(r) \otimes z_r = (b_0(r) z_r, \ldots, b_{P-1}(r) z_r)^\top\).
Position is compressed linearly within each step without shifting the sequence
in time.

Let \(S_{\mathrm{exact}}(z)\) denote the signature of the piecewise-linear path
in \(\R^{RC}\) accumulated from the exact-slot increments \(e_r \otimes z_r\).
Because the accumulated coded path satisfies \(y(t) = (\mathcal A \otimes I_C) x_{\mathrm{exact}}(t)\),
multilinearity of the iterated integrals~\eqref{eq:app-general-signature} makes
each level \(S^k_{\mathcal A}(z) = S^k(y)\) a linear image of the exact-slot
signature:
\begin{equation}\label{eq:offset-projection}
 S^k_{\mathcal A}(z) = (\mathcal A \otimes I_C)^{\otimes k} S^k_{\mathrm{exact}}(z),
 \qquad 1 \leq k \leq m.
\end{equation}

For the normalized Pascal weights \(b_p(r) = \binom{r}{p} / \binom{R-1}{p}\) of
Definition~\ref{def:pascal-lift}, the numerators \(\binom{r}{p} = r(r-1)\cdots(r-p+1)/p!\)
(\(0 \leq p < P\)) form the Newton falling-factorial basis for polynomials of
degree at most \(P-1\). Because a nonzero polynomial of degree below \(P\) can
have at most \(P-1\) roots, every \(P \times P\) minor of the \(P \times R\)
code matrix \(\mathcal A\) is nonzero. At \(P = 1\), \(b_0(r) = 1\) collapses all
slot labels and recovers the ordinary signature of \(z\). For \(1 < P < R\), the
first signature level retains \(P\) polynomial moments per input channel,
\begin{equation}\label{eq:offset-moments}
 s_{c_p} = \sum_{r=0}^{R-1} b_p(r) z_r^c,
\end{equation}
allowing a linear readout to express any degree-\((P-1)\) temporal weighting
\(\sum_{p=0}^{P-1} \beta_{pc} b_p(r)\), while level two adds polynomially
weighted ordered pairs. At the full rank \(P = R\), \(\mathcal A \in \R^{R \times R}\)
is invertible---indeed lower triangular up to column indices, with explicit
binomial inverse
\[
 z_r^c = \sum_{p=r}^{R-1} (-1)^{p-r} \binom{p}{r}\binom{R-1}{p} s_{c_p}
 \qquad (0 \leq r < R)
\]
---so the first-level coded endpoints alone reconstruct the entire sequence
\(z_0, \ldots, z_{R-1}\) linearly.

\subsection{Basis equivalence and the canonical QR orthogonal basis}\label{app:qr}
On the normalized grid \(u_r = r/(R-1)\) (\(R \geq 2\)), each Pascal weight
\(b_p(r)\) is a degree-\(p\) polynomial in \(u_r\) with strictly positive
leading coefficient \((R-1)^p / \bigl(p!\binom{R-1}{p}\bigr)\). For the first
three copies (\(R \geq 3\)),
\[
 b_0(r) = 1, \qquad
 b_1(r) = u_r, \qquad
 b_2(r) = \frac{(R-1)u_r^2 - u_r}{R-2}.
\]
Although \(b_2(r)\) differs from the pure square \(u_r^2\) by a linear term,
the functions \(b_0, \ldots, b_{P-1}\) span the exact same \(P\)-dimensional
polynomial space on \(\{0, \ldots, R-1\}\) as the ascending monomials
\(1, u_r, \ldots, u_r^{P-1}\).

More generally, if \(\widetilde b(r) = L b(r)\) for an invertible transition
matrix \(L \in \R^{P \times P}\), then in copy-major channel order the lifted
path transforms by \(\widetilde y(t) = F y(t)\) with \(F = L \otimes I_C\), and
multilinearity gives
\begin{equation}\label{eq:basis-naturality}
 S^k(\widetilde y) = F^{\otimes k} S^k(y) \quad (1 \leq k \leq m), \qquad
 \widetilde s = F s, \qquad
 \widetilde J = F J F^\top.
\end{equation}
Thus any invertible change of polynomial basis induces an invertible linear
transformation on each complete tensor signature level and on the complete
depth-two logsignature \((s, J)\). Two caveats accompany~\eqref{eq:basis-naturality}:
an invertible basis change need not map a restricted coordinate subset into
itself, and a change of basis can alter the feature Gram matrix governing ridge
regression, even after columnwise standardization
(Subsection~\ref{app:fixed-endpoints}).

Sampled Pascal or monomial columns become increasingly collinear on \([0,1]\) as
\(P\) grows. For the UCI HAR and DSADS representations we therefore orthonormalize
the sampled weights across sequence positions. Let \(A \in \R^{R \times P}\) have
unnormalized binomial entries \(A_{rp} = \binom{r}{p}\) and set
\(D = \operatorname{diag}_p\binom{R-1}{p}\). A thin QR factorization uniquely
writes
\[
 A = Q H, \qquad Q^\top Q = I_P, \qquad H_{pp} > 0 \quad (0 \leq p < P),
\]
where \(Q \in \R^{R \times P}\) has orthonormal columns and \(H \in \R^{P \times P}\)
is upper triangular. Requiring a strictly positive diagonal \(H_{pp} > 0\)
removes the sign ambiguity of the orthonormal columns (the entries
\(q_p(r) = Q_{rp}\) themselves change sign across \(r\) for \(p \geq 1\)). Since
row \(r\) of \(A\) is \(b(r)^\top D\), the orthonormal weight vector satisfies
\(q(r) = (D H^{-1})^\top b(r)\), an explicit invertible basis change of the
form~\eqref{eq:basis-naturality}.

The following is the usual uniqueness of discrete polynomial orthogonalization,
expressed through positive-diagonal QR; see~\cite{discreteqr} for the connection
between discrete orthogonal polynomials and Vandermonde QR factorizations.
\begin{proposition}[Canonical discrete orthogonal basis under positive-diagonal QR]
\label{prop:qr-uniqueness}
For any \(R \geq 2\) and \(1 \leq P \leq R\), let \(U \in \R^{R \times P}\) have
ascending monomial columns \(U_{rp} = u_r^p\) with \(u_r = r/(R-1)\). In exact
arithmetic, the positive-diagonal thin QR factorizations of \(A\), of the
normalized Pascal matrix \(W = A D^{-1}\), and of \(U\) all have the identical
orthonormal factor \(Q\).
\end{proposition}
\begin{proof}
Because \(b_p(r)\) is a degree-\(p\) polynomial in \(u_r\) with positive leading
coefficient, column \(p\) of \(A\) is a linear combination of columns
\(0, \ldots, p\) of \(U\) with a strictly positive coefficient on \(u_r^p\).
Hence \(A = U R_0\) for an upper-triangular matrix \(R_0 \in \R^{P \times P}\)
with positive diagonal entries. Substituting \(A = Q H\) gives
\(U = Q (H R_0^{-1})\) and \(W = Q (H D^{-1})\). Since upper-triangular matrices
with positive diagonals form a group under multiplication and inversion, both
\(H R_0^{-1}\) and \(H D^{-1}\) are upper triangular with positive diagonals.
Uniqueness of the thin QR factorization with positive diagonal entries implies
that all three matrices share the orthonormal factor \(Q\).
\end{proof}

Proposition~\ref{prop:qr-uniqueness} identifies the columns of \(Q\) as the
normalized discrete Legendre polynomials on \(\{0, \ldots, R-1\}\). Applying
Gram--Schmidt to the first two columns \(1\) and \(u_r\) gives their closed form:
\begin{equation}\label{eq:qr-first-two}
 q_0(r) = \frac{1}{\sqrt R}, \qquad
 q_1(r) = \frac{u_r - 1/2}{d_R}, \qquad
 d_R = \sqrt{\sum_{r=0}^{R-1} \Bigl(u_r - \frac{1}{2}\Bigr)^2}
     = \sqrt{\frac{R(R+1)}{12(R-1)}}.
\end{equation}
Under scaled sensor readings \(z_r = a_r / R\), the \(p=0\) endpoint
\(\sum_{r=0}^{R-1} q_0(r) z_r = \bar a / \sqrt R\) preserves the window mean up
to \(R^{-1/2}\). (Ordering channels by sensor axis first and QR copy second, as
in the HAR and DSADS code, simply conjugates \(F\) in~\eqref{eq:basis-naturality}
by a fixed permutation matrix.)

For a unit-sum power spectrum \(p_0, \ldots, p_{K-1}\) on \(K\) equally spaced
frequency bins \(u_k = k/(K-1)\),~\eqref{eq:qr-first-two} also determines the
exact scale of the QR spectral spread coordinate. Because \(J\) is alternating,
subtracting \(\tfrac12 q_0(k)\sqrt K = 1/2\) from \(u_k\) adds a multiple of the
constant column \(q_0\) to \(q_1\) and contributes zero signed area. Hence the
stored HAR spectral area \(\sqrt 2\,J^Q_{01} = (B^Q_{01} - B^Q_{10})/\sqrt 2\) is
an exact positive multiple of the Gini mean difference~\eqref{eq:spectral-separation}:
\begin{equation}\label{eq:spectral-qr-scale}
 \sqrt 2\,J^Q_{01}
   = \frac{\sqrt 2}{\sqrt K\,d_K} J^{(1,u)}_{01}
   = \frac{\sqrt 2}{4\sqrt K\,d_K}\,\E|U - V|,
 \qquad d_K = \sqrt{\frac{K(K+1)}{12(K-1)}},
\end{equation}
where \(U, V\) are independent draws on \(\{u_0, \ldots, u_{K-1}\}\) with
\(\Pr(U = u_k) = p_k\). For an \(N\)-sample time window, physical frequency as a
fraction of the Nyquist rate is \(F = 2k/N = \frac{2(K-1)}{N} U\); thus \(F = U\)
for even \(N = 128\) (\(K = 65\)), and \(F = \frac{126}{127} U\) for
first-difference signals of odd length \(N = 127\) (\(K = 64\)).

\subsection{Shift equivariance and concatenation of locally coded segments}
If positive-diagonal QR supplies the best-conditioned basis on a fixed grid, the
distinctive algebraic virtue of the binomial basis is how it transforms under
index translations: when separately coded adjacent segments are concatenated, the
later segment's local origin must be shifted forward, and in binomial coordinates
that shift is unit lower triangular with integer entries.

Let \(a(r) = \bigl(\binom{r}{0}, \ldots, \binom{r}{P-1}\bigr)^\top \in \R^P\),
where for arbitrary integers \(h \in \Z\) we use the falling-factorial polynomial
\(\binom{h}{p} = h(h-1)\cdots(h-p+1)/p!\) for \(p \geq 1\) and \(\binom{h}{0} = 1\).
(On nonnegative integers \(r \geq 0\), this coincides with setting
\(\binom{r}{p} = 0\) for \(r < p\) in~\eqref{eq:app-pascal} while extending
smoothly to negative shifts.) By Vandermonde's polynomial identity
\(\binom{r+h}{p} = \sum_{q=0}^p \binom{h}{p-q}\binom{r}{q}\), every integer shift
\(h \in \Z\) acts linearly on \(a(r)\):
\begin{equation}\label{eq:offset-shift}
 a(r + h) = T_h a(r), \qquad
 (T_h)_{pq} = \begin{cases}
   \binom{h}{p-q}, & 0 \leq q \leq p < P,\\
   0, & 0 \leq p < q < P.
 \end{cases}
\end{equation}
The matrices \(\{T_h\}_{h \in \Z}\) are unit lower triangular and form an exact
one-parameter abelian group: \(T_0 = I_P\), \(T_h T_k = T_{h+k}\), and
\(T_h^{-1} = T_{-h}\). For \(P = 3\), advancing the step index by \(h = 1\)
applies
\[
 T_1 = \begin{pmatrix} 1 & 0 & 0 \\ 1 & 1 & 0 \\ 0 & 1 & 1 \end{pmatrix},
 \qquad
 T_1 a(2) = \begin{pmatrix} 1 & 0 & 0 \\ 1 & 1 & 0 \\ 0 & 1 & 1 \end{pmatrix}
            \begin{pmatrix} 1 \\ 2 \\ 1 \end{pmatrix}
          = \begin{pmatrix} 1 \\ 3 \\ 3 \end{pmatrix} = a(3).
\]
Because translating the index before coding equals applying \(T_h\) after coding,
the Pascal lift is \emph{shift equivariant}. In normalized Pascal coordinates
\(b(r) = D^{-1} a(r)\) with \(D = \operatorname{diag}_p\binom{R-1}{p}\), or after
a further invertible basis map \(L\) such as QR, the shift operator conjugates to
\(D^{-1} T_h D\) or \(L D^{-1} T_h D L^{-1}\).

Combining shift equivariance with Chen's identity~\cite{chen} yields a
concatenation rule for locally coded segments without revisiting raw increments.
For any segment \(z = (z_0, \ldots, z_{\ell-1})\) of
length \(\ell \geq 1\) coded from its own local origin \(j = 0, \ldots, \ell-1\)
in unnormalized Pascal coordinates, its truncated depth-\(m\) tensor signature is
\begin{equation}\label{eq:reset-signature}
 G(z) = \prod_{j=0}^{\ell-1} \exp_\otimes\bigl(a(j) \otimes z_j\bigr).
\end{equation}
Let \(\Theta_h\) denote the truncated tensor-algebra automorphism that applies
\(T_h \otimes I_C\) to each tensor factor:
\(\Theta_h|_{(\R^{PC})^{\otimes k}} = (T_h \otimes I_C)^{\otimes k}\). Since
\((T_\ell \otimes I_C)\bigl(a(j) \otimes z_j\bigr) = a(\ell + j) \otimes z_j\)
and \(\exp_\otimes\) commutes with linear maps, concatenating \(z\) with a second
locally coded segment \(w = (w_0, \ldots, w_{q-1})\) satisfies the
\emph{twisted Chen identity}
\begin{equation}\label{eq:reset-chen}
 G(z \mathbin\Vert w) = G(z)\,\Theta_\ell\bigl(G(w)\bigr).
\end{equation}
Associativity across multiple blocks follows immediately from the group law
\(\Theta_\ell \Theta_q = \Theta_{\ell+q}\) and associativity of the truncated
tensor product. If segments of different lengths \(\ell\) were normalized
independently by their own denominators \(\binom{\ell-1}{q}\) (using
\(K = \min(P, \ell)\) local copies so even a singleton step \(\ell = 1\) avoids
division by zero), applying the \(P \times K\) rebasing map
\(H_{pq} = \binom{h}{p-q}\binom{\ell-1}{q}/\binom{R-1}{p}\) (\(q \leq p\)) to a
block starting at global index \(h\) embeds its local signature directly into the
global normalized signature \(S_{\leq m}(y)\). Under uniform sampling, integer
index shifts correspond to physical time translations.

\begin{remark}[Online sliding-window updates and comparison with SigAdapt]
\label{app:sigadapt}
Equation~\eqref{eq:reset-chen} also clarifies how representation choices differ
from online readout adaptation, such as SigAdapt~\cite{sigadapt} for sliding-window
time-series forecasting. For each look-back length, SigAdapt computes ordinary
tensor signatures (at depths four to six in its reported experiments) of
time- and basepoint-augmented history and covariate paths, fits a joint ridge
model, updates multipliers of its per-level signature contributions online via a
Kalman filter, and aggregates across window lengths. Its sliding-window signature
updates drop the oldest step and append the newest step via Chen's identity and basepoint
correction. Performing an analogous sliding update on locally indexed Pascal
signatures additionally applies \(\Theta_{-1}\) after left-multiplying by
\(\exp_\otimes\bigl(-a(0)\otimes z_0\bigr)\) so all remaining step indices shift
down by one before the new step is multiplied on the right at \(a(R-1)\).
Polynomial coding and online state-space adaptation address orthogonal parts of a
modelling pipeline and have not been combined here; moreover, algebraic update
identities for raw signature blocks do not determine end-to-end cost when a
representation also recomputes windowed Fourier transforms or lagged geometric
streams (Sections~\ref{app:har} and~\ref{app:compute}).
\end{remark}

\subsection{Nilpotent generator and Volterra signature specialization}
\label{app:nilpotent}
Because \(\{T_h\}_{h \in \Z}\) is a unipotent lower-triangular group, the Pascal
weights arise from a nilpotent matrix exponential, and the entire coded signature
embeds into the Volterra signature framework~\cite{volterra}. It is instructive
to derive both descriptions and keep their scopes distinct: the first is an
identity for single-step weights, whereas the second identifies iterated
integrals across all signature depths.

\begin{proposition}[Nilpotent generator of Pascal weights]\label{prop:nilpotent}
Let \(N \in \R^{P \times P}\) have ones on its first subdiagonal
(\(N_{p, p-1} = 1\) for \(1 \leq p < P\)) and zeros elsewhere, let \(e_0\) be the
first standard basis vector of \(\R^P\), and define the finite matrix logarithm
\[
 H = \log(I_P + N) = \sum_{j=1}^{P-1} \frac{(-1)^{j+1}}{j} N^j.
\]
Then \(T_h = e^{hH}\) and \(a(r) = e^{rH} e_0\) for every integer \(h, r \in \Z\).
In particular, relative to any fixed terminal index \(T \in \Z\),
\begin{equation}\label{eq:nilpotent-endpoint}
 e^{-(T - r) H} e_0 = e^{-T H} a(r).
\end{equation}
\end{proposition}
\begin{proof}
Because \(N^P = 0\), the power series for \(\log(I_P + N)\) and \(\exp(hH)\)
truncate at degree \(P - 1\), and formal power-series composition holds exactly:
\(e^H = I_P + N\) and \(e^{hH} = (I_P + N)^h\). By the generalized binomial
theorem for nilpotent perturbations,
\[
 (I_P + N)^h = \sum_{k=0}^{P-1} \binom{h}{k} N^k \qquad (h \in \Z).
\]
Since \((N^k)_{pq} = \mathbf 1_{\{p - q = k\}}\), entry \((p, q)\) of
\((I_P + N)^h\) equals \(\binom{h}{p-q}\) for \(q \leq p\) and \(0\) otherwise,
which is \(T_h\) in~\eqref{eq:offset-shift}. Evaluating \(T_r e_0\) selects column
\(q = 0\) and gives \(a(r)\); left-multiplying by \(e^{-TH}\) yields~\eqref{eq:nilpotent-endpoint}.
\end{proof}

Polynomial memory from nilpotent generators is treated within exponentially
weighted signatures (EWS) in~\cite[Prop.~2]{ews}. EWS also identifies its
fixed-horizon iterated integrals with the ordinary signature of a reweighted
path~\cite[Prop.~3]{ews}. For the present discrete lift, holding the clock at
index \(r\) during signal increment \(z_r\) makes the terminal-weighted increments
in~\eqref{eq:nilpotent-endpoint} a fixed linear transform of the Pascal increments;
the corresponding signatures then transform level by level. This comparison
uses that clock convention, rather than simultaneous linear clock and signal
interpolation.

To place the coded signature within the Volterra signature~\cite{volterra} at
every truncation level, recall that a Volterra signature replaces ordinary path
differentials by integrals against a measurable matrix-valued kernel \(K(t, s)\)
satisfying \(\sup_{t \in [0, R]} \int_0^t \|K(t, s)\|\,ds < \infty\). Let
\(x:[0, R]\to\R^C\) be the piecewise-linear path with increments \(z_r\) on
\([r, r+1]\), and define the piecewise-constant, position-only kernel
\[
 A(s) = b(r) \otimes I_C \quad (r \leq s < r + 1,\ 0 \leq r < R), \qquad
 K(t, s) = A(s) \quad (0 \leq s \leq t \leq R),
\]
with \(A(R) = b(R-1) \otimes I_C\). Because \(K\) is bounded on \([0, R]^2\), it
is Volterra admissible, and integrating \(A(s)\,dx_s\) on \([0, t]\) produces the
exact piecewise-linear lifted path \(y_t = \int_0^t A(s)\,dx_s\) of
Definition~\ref{def:pascal-lift}. At signature level \(k \geq 1\), the Volterra
iterated integral~\cite[Def.~2.14]{volterra} over the simplex
\(0 < r_1 < \cdots < r_k < t\) with terminal horizon \(r_{k+1} = \tau \geq t\)
forms the tensor product of factors \(K(r_{\ell+1}, r_\ell)\,dx_{r_\ell}\).
Because \(K(r_{\ell+1}, r_\ell) = A(r_\ell)\) is independent of its first time
argument \(r_{\ell+1}\), the \(\ell\)th differential simplifies to
\(A(r_\ell)\,dx_{r_\ell} = dy_{r_\ell}\), yielding exact equality at all levels:
\[
 \VSig(x; K)^{\tau}_{0, t} = \Sig(y)_{0, t} \qquad (0 \leq t \leq \tau \leq R).
\]
Subsequent rescaling or QR orthonormalization applies a fixed linear map
\(F^{\otimes k}\) to each level. Because \(K(t, s)\) depends only on \(s\), this
specialization preserves the classical Chen concatenation and shuffle identities
of ordinary signatures.

By contrast, while~\eqref{eq:nilpotent-endpoint} shows that weighting discrete
increments by the terminal-lag exponential \(e^{-(T-r)H} e_0\) differs from
\(a(r) \otimes z_r\) only by the fixed invertible transformation
\(e^{-TH} \otimes I_C\), that relation does \emph{not} identify higher levels of
the Pascal-coded signature with a convolution Volterra signature
\(K(t, s) = \widetilde K(t - s)\). In a genuine two-time convolution kernel, the
\(\ell\)th Volterra factor \(K(r_{\ell+1}, r_\ell)\) weights step \(r_\ell\) by
its lag relative to the \emph{next} integration variable \(r_{\ell+1}\), rather
than referencing all factors to a single terminal endpoint \(T\). This distinction
between fixed-terminal EWS and convolution Volterra weighting is also discussed
in~\cite[App.~D.2]{ews}.
